\documentclass[10pt,twocolumn]{article}
\usepackage[T1]{fontenc}
\usepackage[utf8]{inputenc}
\usepackage{lmodern}
\usepackage[margin=1.8cm,top=2.4cm,bottom=2cm,columnsep=0.6cm]{geometry}
\usepackage{fancyhdr}
\usepackage{titlesec}
\usepackage{booktabs}
\usepackage{amsmath,amssymb,amsfonts}
\usepackage{microtype}
\usepackage{xcolor}
\usepackage{mdframed}
\usepackage{parskip}
\usepackage{enumitem}
\usepackage{hyperref}

\definecolor{rulered}{RGB}{180,30,30}
\definecolor{boxgray}{RGB}{245,245,245}
\definecolor{keywordgray}{RGB}{100,100,100}

\pagestyle{fancy}
\fancyhf{}
\fancyhead[L]{\textsc{\small Claude Mag}}
\fancyhead[C]{\small\textit{Machine Learning Research Digest}}
\fancyhead[R]{\small 2026-10-05}
\fancyfoot[C]{\small\thepage}
\renewcommand{\headrulewidth}{0.8pt}

\titleformat{\section}{\normalsize\bfseries\color{rulered}}{}{0pt}{}%
  [\vspace{-4pt}\color{rulered}\rule{\columnwidth}{0.4pt}]
\titlespacing{\section}{0pt}{8pt}{4pt}

\hypersetup{colorlinks=true,urlcolor=rulered,linkcolor=black}

\begin{document}
\setlength{\parindent}{0pt}
\setlength{\parskip}{4pt}

{\small\color{keywordgray}\textbf{Keywords:}
  mathematical reasoning \textbullet{}
  LLM interpretability \textbullet{}
  structural understanding \textbullet{}
  primitive-guided post-training}\par\vspace{4pt}

{\LARGE\bfseries\raggedright
  The Missing Primitive: Diagnosing and Repairing Mathematical Reasoning in Large Language Models}\par\vspace{4pt}

{\large\itshape\raggedright
  A New Benchmark Reveals that Discovery---Not Execution---Is the Universal Bottleneck
  to LLM Mathematical Reasoning, and That Targeting It Yields Consistent Gains}\par\vspace{6pt}

{\small\textbf{By} Shuo Xing, Zilin Dai, Chengyuan Qian, Fangzhou Lin, Wenjing Chen,
  Ping He, Pan Lu, Alvaro Velasquez, Mohit Bansal \& Zhengzhong Tu
  (Texas A\&M University; UNC Chapel Hill; UCLA; University of Colorado Denver)}\par\vspace{2pt}

{\small\color{keywordgray}arXiv:2610.02191 \textbullet{} October 2026}
\par\vspace{2pt}\noindent{\color{rulered}\rule{\columnwidth}{0.6pt}}\vspace{6pt}

LLMs that routinely solve competition-level mathematics often lack the structural understanding
underlying their solutions: given the correct foundational idea---a \emph{Mathematical Primitive}---the
same models execute near-perfectly, but independently discovering that primitive is the dominant
bottleneck across all frontier architectures tested. Primitive-guided post-training closes the gap.

\begin{mdframed}[backgroundcolor=boxgray,linecolor=rulered,linewidth=1pt,
    innerleftmargin=6pt,innerrightmargin=6pt,
    innertopmargin=6pt,innerbottommargin=6pt]
\textbf{\small AT A GLANCE}\par\vspace{4pt}
\begin{description}[leftmargin=0pt,labelsep=4pt,itemsep=2pt,parsep=0pt,
    font=\small\bfseries]
  \item[Problem:] LLM mathematical reasoning accuracy does not reveal whether models
    possess the structural understanding behind their solutions.
  \item[Why it matters:] Models with identical benchmark scores may rely on brittle
    heuristics; diagnosing structural gaps enables targeted improvement.
  \item[Method:] Mathematical Primitive concept plus the Prim benchmark spanning
    Discovery, Generation, Digestion, and Execution dimensions.
  \item[Result:] Discovery lags Execution by up to 30 pp; primitive-guided post-training
    consistently raises both Discovery and overall task accuracy.
  \item[Evidence:] Evaluated across multiple frontier LLMs on Prim and held-out
    downstream math benchmarks.
\end{description}
\end{mdframed}

\section{The Big Picture}

Large language models have crossed remarkable thresholds in mathematical reasoning: the best
contemporary systems solve a substantial fraction of AMC and AIME problems and approach expert
performance on graduate-level benchmarks. Yet this aggregate success conceals a potentially
critical gap: does solving a problem indicate that the model grasped the foundational idea
organizing the solution, or merely that it reproduced a pattern from training data?

This paper takes a first step toward answering that question by introducing the concept of
the \emph{Mathematical Primitive}---a concise, non-procedural description of the core insight
that, once identified, makes the detailed execution steps straightforward. For example,
a combinatorics problem's primitive might be ``map to a stars-and-bars counting argument''
rather than the specific arithmetic that follows from it.

\section{Why Existing Approaches Fall Short}

\textbf{Benchmark saturation:} State-of-the-art LLMs now score near ceiling on MATH and
AMC; high scores mask vastly different capability profiles.

\textbf{No structural audit:} Existing benchmarks measure correct final answers, not whether
the model identified the organizing idea. Two models with identical accuracy may differ
dramatically in structural understanding.

\textbf{Post-training targeting:} Standard chain-of-thought fine-tuning and process reward
models treat all reasoning steps equally; they do not specifically target primitive
discovery where diagnosis reveals the greatest weakness.

\section{Core Mechanism}

\textbf{Mathematical Primitive.} A primitive is defined as the single, concise, non-procedural
representation of the foundational idea that organizes a solution. It is shorter than a proof
sketch, more specific than a problem category, and must identify the key conceptual move.

\textbf{Prim Benchmark.} The paper constructs Prim with four complementary dimensions:
\begin{itemize}[itemsep=2pt,parsep=0pt]
  \item \emph{Discovery:} Given the problem, identify the primitive.
  \item \emph{Generation:} Given the primitive, produce a valid solution.
  \item \emph{Digestion:} Given a complete solution, extract the primitive it relies on.
  \item \emph{Execution:} Given both problem and primitive, carry out the solution steps.
\end{itemize}

\section{Technical Formulation}

Let $\mathcal{P}$ be the space of primitives and $\mathcal{Q}$ the space of problems.
Define four binary capability indicators for a model $\mathcal{M}$ on problem $q$ with
ground-truth primitive $p^*(q)$:
\[
\text{Disc}(q) = \mathbf{1}[\hat{p}(q) = p^*(q)]
\]
\[
\text{Exec}(q,p^*) = \mathbf{1}[\hat{s}(q,p^*(q))\text{ is correct}]
\]
The central finding is:
\[
\mathbb{E}_q\bigl[\text{Exec}(q,p^*(q))\bigr] \gg \mathbb{E}_q\bigl[\text{Disc}(q)\bigr]
\]
for all tested models, establishing that \textbf{Discovery is the dominant bottleneck}.
Primitive-guided post-training minimizes a combined loss:
\[
\mathcal{L} = \mathcal{L}_{\text{solve}} + \lambda \mathcal{L}_{\text{discover}}
\]
where $\mathcal{L}_{\text{discover}}$ is a cross-entropy term on primitive Discovery, and
$\lambda$ is a balancing hyperparameter.

\section{Learning or Inference Procedure}

\textbf{Evaluation:}
\begin{enumerate}[itemsep=2pt,parsep=0pt]
  \item Run each Prim dimension independently on the same model.
  \item Compute per-dimension accuracy and capability profile.
  \item Compare profiles across models with similar aggregate task accuracy.
\end{enumerate}

\textbf{Post-training:}
\begin{enumerate}[itemsep=2pt,parsep=0pt]
  \item Annotate training problems with ground-truth or LLM-generated primitives.
  \item Add a primitive Discovery objective alongside the standard solution objective.
  \item Fine-tune with combined loss $\mathcal{L}$; evaluate on held-out Prim dimensions
    and downstream benchmarks.
\end{enumerate}

\section{What the Guarantee Says}

No formal convergence bound is given; the Discovery bottleneck is an empirical finding.
The central claim is: across all tested frontier LLMs, Execution given a correct primitive
substantially exceeds unconditioned Discovery, establishing the latter as universally limiting.
Post-training improvement transfers to held-out benchmarks not seen during fine-tuning.

\section{Experimental Findings}

\begin{itemize}[itemsep=2pt,parsep=0pt]
  \item \textbf{Discovery gap:} Execution given the correct primitive exceeds unconditioned
    Discovery by up to 30 percentage points across all frontier LLMs.
  \item \textbf{Hidden differences:} Two models with identical MATH accuracy can differ
    by over 20 pp on the Discovery dimension alone.
  \item \textbf{Latent execution capacity:} Providing correct primitives to models that
    cannot discover them independently unlocks near-full solution accuracy.
  \item \textbf{Post-training gains:} Primitive-guided fine-tuning consistently improves
    Discovery accuracy and overall task performance, with gains transferring to held-out sets.
\end{itemize}

\section{Ablations and Interpretation}

\textbf{Primitive granularity:} Coarser primitives yield lower Execution improvement;
fine-grained primitives that specify the exact conceptual move are most effective.

\textbf{Imperfect primitives:} Using LLM-generated rather than ground-truth primitives
reduces but does not eliminate the Execution advantage---even imperfect discovery guidance helps.

\textbf{Digestion vs.\ Discovery:} High Digestion accuracy does not predict high Discovery
accuracy, establishing that recognizing a primitive in a given solution and independently
finding it from a problem statement are distinct capabilities.

\section*{Reference}

Shuo Xing et al.
\textbf{The Missing Primitive: Diagnosing and Repairing Mathematical Reasoning in
Large Language Models}.
arXiv:2610.02191, October 2026.\\
\url{https://arxiv.org/abs/2610.02191}

\end{document}
